\section*{Chapitre \ref{ch:b3:developmental-genetics} --- Génétique du développement et évolution des formes}

\begin{solution}{exo:b3:developmental-genetics:1}
Effet maternel~: \emph{bicoid}, les embryons de mères mutantes n'ont ni
tête ni thorax. Gap~: \emph{Krüppel}, un bloc contigu de segments
manquant. Règle de parité~: \emph{fushi tarazu} (ou \emph{even-skipped}),
un segment sur deux manquant. Polarité segmentaire~: \emph{hedgehog} (ou
\emph{engrailed}), une partie de chaque segment remplacée par l'image en
miroir du reste.
\end{solution}

\begin{solution}{exo:b3:developmental-genetics:2}
Un gène dont le produit est déposé dans l'œuf par la mère, de sorte que
le phénotype de l'embryon suit le génotype de celle-ci. Ses embryons
n'ont ni tête ni thorax et portent des structures postérieures aux deux
bouts~; ils meurent. L'allèle sauvage paternel est dans le zygote, mais
l'ARNm de \emph{bicoid} est fabriqué pendant l'ovogenèse par les
cellules nourricières de la mère et ancré au pôle avant la fécondation~;
la copie du zygote ne serait transcrite que des heures plus tard, une
fois passé le temps du gradient.
\end{solution}

\begin{solution}{exo:b3:developmental-genetics:3}
\omterm{def:b3:developmental-genetics:hox}{Colinéarité}~: l'ordre des \omterm{def:b3:developmental-genetics:hox}{gènes Hox} le long du chromosome correspond à
l'ordre de leurs domaines d'expression le long du corps. Prévalence
postérieure~: là où plusieurs sont exprimés, le plus postérieur fixe
l'identité. Une souris dépourvue de \emph{Hoxc8} montre une
transformation antérieure --- sa première vertèbre lombaire prend
l'identité d'une thoracique et porte une côte. \emph{Antennapedia} dans
la tête change les antennes en pattes, l'identité du deuxième segment
thoracique.
\end{solution}

\begin{solution}{exo:b3:developmental-genetics:4}
Une lèvre dorsale greffée a induit un second axe complet fait surtout de
cellules de l'hôte~: le greffon instruisait ses voisines --- un
\omterm{def:b3:developmental-genetics:organiser}{organisateur}. Ses molécules sont des antagonistes sécrétés de BMP
(Chordine, Noggine, Follistatine). L'ectoderme devient neural quand la
signalisation BMP est absente et épiderme quand BMP est présent, de
sorte que le neural est ce que fait l'ectoderme quand on ne lui dit
rien~: l'\omterm{def:b3:developmental-genetics:organiser}{organisateur} agit en faisant taire une instruction.
\end{solution}

\begin{solution}{exo:b3:developmental-genetics:5}
$k = \ln 2/(35\times 60) = \qty{3.3e-4}{s^{-1}}$~; $\lambda = \sqrt{4/
3.3\times 10^{-4}} = \qty{110}{\micro m}$~; rapport $\mathrm{e}^{250/110}
= \mathrm{e}^{2.27} = 9.7$.
\end{solution}

\begin{solution}{exo:b3:developmental-genetics:6}
Quatre copies doublent $c_{0}$~: la frontière recule de $\lambda\ln 2 =
\qty{69}{\micro m}$ jusqu'à \qty{309}{\micro m}~; une copie le divise par
deux~: elle avance de \qty{69}{\micro m} jusqu'à \qty{171}{\micro m}.
Chaque décalage fait \qty{14}{\%} de la longueur de l'œuf --- plus que
ce qu'on observe, ce qui est déjà un indice de mécanismes qui amortissent
la dépendance à la dose.
\end{solution}

\begin{solution}{exo:b3:developmental-genetics:7}
$\delta x = \lambda\,\delta c/c$~: pour \qty{8}{\micro m}, $\delta c/c =
8/100 = \qty{8}{\%}$. Un comptage de Poisson $N$ a une erreur relative
$1/\sqrt{N}$, donc il faut compter $N = 1/0.08^{2} \approx 160$
molécules --- quelques pour cent des molécules présentes dans un noyau,
ou un comptage unique intégré dans le temps.
\end{solution}

\begin{solution}{exo:b3:developmental-genetics:8}
$k = \ln 2/35 = \qty{0.0198}{min^{-1}}$~: $1 - \mathrm{e}^{-kt}$ donne
$0.70$ à \qty{60}{min}, $0.83$ à \qty{90}{min}, $0.91$ à
\qty{120}{min}, $0.95$ à \qty{150}{min}. Le gradient monte encore de
plusieurs pour cent pendant qu'on le lit, ce qui décalerait une frontière
de $\lambda\ln(1/0.9) \approx \qty{10}{\micro m}$, environ un noyau, sur
la fenêtre de lecture~: l'embryon doit soit lire à un instant fixe, soit
employer une lecture insensible au niveau d'ensemble.
\end{solution}

\begin{solution}{exo:b3:developmental-genetics:9}
Une seconde ZAP donne une duplication en miroir, 4-3-2-2-3-4~: les
cellules antérieures voient maintenant beaucoup de Shh et deviennent des
doigts postérieurs. Une bille qui libère peu de Shh donne une
duplication partielle --- un doigt 2 antérieur en plus, soit 2-2-3-4 ---
parce que les cellules antérieures ne voient que la faible concentration
qui spécifie le doigt 2.
\end{solution}

\begin{solution}{exo:b3:developmental-genetics:10}
Solution particulière $s/k$~; solutions homogènes $\cosh(x/\lambda)$,
$\sinh(x/\lambda)$~; l'absence de flux en $0$ élimine le $\sinh$~;
$c(L) = 0$ fixe l'amplitude~:
$c(x) = (s/k)\bigl[1 - \cosh(x/\lambda)/\cosh(L/\lambda)\bigr]$.
Le profil est un plateau près de $x = 0$ qui tombe à zéro dans une
couche limite de largeur $\lambda$ au puits --- non pas une exponentielle
issue d'un point, mais un champ plat avec un bord~; la concentration la
plus haute est la plus éloignée du puits. Pour $L \ll \lambda$,
$\cosh u \approx 1 +
u^{2}/2$ et $c \approx s(L^{2} - x^{2})/(2D)$~: une parabole indépendante
de $k$, puisque les molécules atteignent le puits avant de pouvoir se
dégrader.
\end{solution}

\begin{solution}{exo:b3:developmental-genetics:11}
$x_{1} = \lambda\ln(1/0.3) = 1.20\lambda$ et $x_{2} = \lambda\ln(1/0.08)
= 2.53\lambda$~: la première région fait $1.20\lambda$ de large, la
seconde $1.32\lambda$, toutes deux indépendantes de $c_{0}$, qui ne fait
que les décaler ensemble. La troisième région, $L - 2.53\lambda$, absorbe
toute la variation de longueur de l'œuf~: dans un œuf de
\qty{450}{\micro m} avec $\lambda = \qty{100}{\micro m}$ l'abdomen fait
\qty{197}{\micro m}, dans un œuf de \qty{550}{\micro m}
\qty{297}{\micro m} --- le motif n'est pas mis à l'échelle. Mécanismes~:
un second gradient venu du pôle postérieur (Nanos, Caudal, ou le système
terminal) lu de concert avec Bicoid, de sorte que les positions soient
fixées par le rapport~; ou une molécule «~extenseur~» faite là où le
morphogène est faible et qui allonge $\lambda$ jusqu'à ce que le gradient
remplisse l'œuf.
\end{solution}

\begin{solution}{exo:b3:developmental-genetics:12}
Un changement codant altère la protéine partout où elle agit~:
\emph{Pitx1} bâtit aussi l'\omterm{def:b3:endocrinology:axes}{hypophyse} et la mâchoire, \emph{\omterm{def:b3:developmental-genetics:organiser}{sonic
hedgehog}} structure aussi le tube neural, l'intestin et les dents, de
sorte qu'une protéine brisée est létale ou pléiotropiquement
invalidante. Un activateur commande l'expression dans un tissu à un
moment~; le supprimer retire la structure que ce tissu aurait bâtie et
laisse intacts tous les autres usages de la protéine. Les activateurs
sont modulaires et souvent partiellement redondants, de sorte que les
étapes intermédiaires sont viables, et la même route --- un
interrupteur perdu --- a été prise indépendamment dans de nombreux lacs
à épinoches.
\end{solution}

\begin{solution}{pb:b3:developmental-genetics:1}
\textbf{1.} $k = \ln 2/(35\times 60) = \qty{3.3e-4}{s^{-1}}$~; $1/k =
\qty{3030}{s} = \qty{50}{min}$.
\textbf{2.} $\lambda = \sqrt{4/3.3\times 10^{-4}} = \qty{110}{\micro m}$~;
l'œuf fait $4.5\lambda$.
\textbf{3.} $\mathrm{e}^{4.5} \approx 90$ d'un pôle à l'autre~;
$\mathrm{e}^{8.5/110} = 1.08$, soit \qty{8}{\%} de changement par noyau.
\textbf{4.} $J = 50\times\sqrt{4\times 3.3\times 10^{-4}} = 50\times
0.036 = \qty{1.8}{nM\,\micro m\,s^{-1}}$. Un nM fait $0.6$ molécule
par \unit{\micro m^3}, donc environ $1.1$ molécule par \unit{\micro m^2}
et par seconde.
\textbf{5.} Volume nucléaire $\tfrac{4}{3}\pi\times 27 = \qty{113}{\micro
m^3}$. À \qty{50}{nM}, $30$ molécules par \unit{\micro m^3}~:
\num{3400} molécules. Au milieu de l'œuf $c = 50\,\mathrm{e}^{-2.27} =
\qty{5.2}{nM}$~: environ \num{350} molécules.
\textbf{6.} $kt = 1.19$ à \qty{60}{min}~: $0.70$~; $kt = 2.38$ à
\qty{120}{min}~: $0.91$. Le gradient est à \qty{10}{\%} de l'état
stationnaire quand il est lu --- il monte encore d'une quantité qui
déplacerait une frontière de $\lambda\ln(1/0.91) \approx \qty{10}{\micro m}$,
un noyau.
\textbf{7.} $x_{1} = 110\ln(1/0.3) = \qty{132}{\micro m}$
(\qty{26}{\%})~; $x_{2} = 110\ln(1/0.08) = \qty{278}{\micro m}$
(\qty{56}{\%}).
\textbf{8.} Quatre copies~: $c_{0} = \qty{100}{nM}$~; les deux frontières
reculent de $\lambda\ln 2 = \qty{76}{\micro m}$, à $208$ et
\qty{354}{\micro m}. Six copies~: $\lambda\ln 3 = \qty{121}{\micro m}$,
à $253$ et \qty{399}{\micro m}.
\textbf{9.} Une copie~: elles avancent de \qty{76}{\micro m}, à $56$ et
\qty{202}{\micro m}. La région de la tête a rétréci de $132$ à
\qty{56}{\micro m}~; le thorax a gardé sa largeur (\qty{146}{\micro m}) et
s'est déplacé vers l'avant~; l'abdomen est passé de $222$ à \qty{298}{\micro m}.
\textbf{10.} $\delta x = 110\times 0.1 = \qty{11}{\micro m}$, soit $1.3$
espacement nucléaire.
\textbf{11.} Un noyau~: $\delta c/c = 8.5/110 = \qty{7.7}{\%}$. À partir
de \qty{10}{\%}, $n = (10/7.7)^{2} \approx 1.7$~: deux lectures
indépendantes (ou la moyenne de deux noyaux voisins) suffisent.
\textbf{12.} $1/\sqrt{350} = \qty{5.3}{\%}$~: un comptage instantané
unique est déjà plus précis que les \qty{10}{\%} supposés, de sorte que
le bruit observé n'est pas seulement un bruit de comptage mais vient de
la cinétique de liaison et de la machinerie de lecture --- et moyenner
dans le temps fait encore mieux.
\textbf{13.} $278/450 = \qty{62}{\%}$ et $278/550 = \qty{51}{\%}$ de la
longueur de l'œuf~: un écart de \qty{11}{\%}, six noyaux. Le motif n'est
pas mis à l'échelle avec un $\lambda$ fixe.
\textbf{14.} Si $D \propto L^{2}$ à $k$ fixé, $\lambda \propto L$ et
$x_{T}/L = (\lambda/L)\ln(c_{0}/T)$ est le même dans tout œuf~: mise à
l'échelle parfaite. Mais $D$ est une propriété de la molécule et du
cytoplasme, non de la taille de l'œuf, de sorte que cela est
invraisemblable tel quel~; la mise à l'échelle des vrais embryons est
partielle et vient d'autres lectures.
\textbf{15.} $c^{5}/(c^{5} + T^{5}) = 0.1$ pour $c/T = 9^{-1/5} = 0.64$ et
$0.9$ pour $c/T = 9^{1/5} = 1.55$~: un facteur $2.4$ en concentration,
$\Delta x = 110\ln 2.4 = \qty{97}{\micro m}$ --- onze noyaux, bien plus
large que la frontière observée de deux ou trois~; la seule lecture du
gradient n'est pas assez nette.
\textbf{16.} La raideur convertit une concentration donnée en allumé ou
éteint de façon plus tranchée, de sorte que la région de transition est
plus étroite~; mais une erreur de \qty{10}{\%} sur la concentration
déplace toujours de $\lambda\,\delta c/c$ le point où le seuil est
franchi~: la réponse peut être une marche parfaite et la marche tomber
quand même au mauvais endroit.
\textbf{17.} Si les deux étaient allumés, chacun réprimerait l'autre, de
sorte qu'au moins l'un serait éteint~; avec une répression coopérative
les seuls états stables sont l'un allumé et l'autre éteint, et une
frontière entre eux est un interrupteur, non une pente.
\textbf{18.} $\mathrm{d}\ln\lambda = \tfrac{1}{2}(0.02 - 0.06) = -0.02$
par degré~: $\lambda$ baisse de \qty{2}{\%} par degré, \qty{10}{\%} pour
\qty{5}{\celsius}, jusqu'à \qty{99}{\micro m}~; la frontière de
\emph{hunchback} passe de $278$ à \qty{250}{\micro m}, trois noyaux vers
l'avant. Les mouches se développent entre \qty{18}{} et
\qty{29}{\celsius} avec des proportions normales, donc quelque chose
compense --- une source ou une lecture qui se décale en sens inverse
avec la température.
\textbf{19.} $2^{8} = 256$ codes~; quatorze employés. Avec la prévalence
postérieure seul compte le gène allumé le plus postérieur, donc au plus
neuf identités distinctes (huit gènes, ou aucun)~: le code est un
classement, non un mot binaire.
\textbf{20.} Le troisième segment thoracique devient un second~:
\emph{Ubx} réprimait le programme de l'aile qu'un segment thoracique
exécute par défaut, et retirer le répresseur libère l'état ancestral ---
quatre ailes, comme chez les ancêtres à quatre ailes de la mouche.
\textbf{21.} La région cervicale de l'oie va jusqu'à la vertèbre 22~: un
long cou de nombreuses vertèbres, contre sept chez la souris. Le gène
est le même~; la limite antérieure de son expression a reculé --- un
changement régulateur du lieu d'allumage d'un gène Hox.
\textbf{22.} Chaque marge est maintenant une source~; Shh est le plus
élevé aux deux bords (doigt 4), tombe vers le centre (3, puis 2), et
n'atteint jamais le faible niveau qui spécifie le doigt 1, de sorte que
le milieu tient deux doigts 2 dos à dos~: 4-3-2-2-3-4.
\textbf{23.} \emph{Pax6} est un régulateur maître conservé qui peut
déclencher n'importe quel programme d'œil que porte l'hôte~: la mouche a
bâti un œil de mouche, donc le gène de souris ne porte aucune
information d'œil de souris, seulement l'ordre «~bâtis un œil ici~».
Cela montre l'\omterm{def:b3:developmental-genetics:evodevo}{homologie profonde} du réseau de gènes, non l'homologie des
organes~: l'œil composé et l'œil à chambre sont apparus séparément à
partir d'un ancêtre qui avait une plage de photorécepteurs à
\emph{Pax6}.
\textbf{24.} Cinq doigts en travers de \qty{2}{mm} font une longueur
d'onde de \qty{0.4}{mm}~; six en demandent \qty{0.33}{mm}, un
raccourcissement d'un sixième. Puisque
$\lambda \propto \sqrt{D_{\text{inh}}/k_{\text{inh}}}$, il faut élever la
dégradation de l'inhibiteur de \qty{44}{\%} ou abaisser sa diffusion de
\qty{31}{\%}~; dans la plaque de la main, réduire la dose des \omterm{def:b3:developmental-genetics:hox}{gènes Hox}
postérieurs y parvient, et les souris mutantes ont six, sept doigts fins
ou davantage.
\textbf{25.} $\lambda \approx \qty{110}{\micro m}$~; un doublement de la
source décale chaque frontière de \qty{76}{\micro m}~; un bruit de
\qty{10}{\%} sur la concentration fait \qty{11}{\micro m}, environ un
noyau, d'erreur de position.
\end{solution}
